%\clearpage

\section{Summary of changes included in version 7 of T2K-TN-141 (released on June 18, 2013)}

This section is designed to allow readers to keep track of the changes between versions of T2K-TN-141.
Table, figure, equation, footnote or citation numbers shown in black refer to version of 6,
{\color{red}whereas the numbers shown in red refer to version of 7}.
Both numbers are shown where possible.

\begin{itemize}

\item
  At the request of Nakaya-san, we now quote the
  68\% error on $\sin^2 2\theta_{23}$, and the 90\% C.L. interval for
  $|\Delta m^2_{32}|$ in Section 5.4 {\color{red}(5.4)}.
  This is in addition to the 90\% C.L. limit for $\sin^2 2\theta_{23}$ and
  the 68\% error for $|\Delta m^2_{32}|$ already presented.

\item
  At the request of the paper committee, we now quote the edges
  of the '2+N' Feldman Cousins contour in Section 5.5 {\color{red}(5.5)}.
  The boundary of the 68\% '2+N' contour includes the points
  ($\sin^2(2\theta_{23})$,$|\Delta m^2_{32}|$) =
  (1.000,2.64$\times10^{-3}$ eV$^2$/c$^4$),
  (0.964,2.50$\times10^{-3}$ eV$^2$/c$^4$) and
  (1.000,2.26$\times10^{-3}$ eV$^2$/c$^4$).
  The boundary of the 90\% '2+N' contour includes the points
  ($\sin^2(2\theta_{23})$,$|\Delta m^2_{32}|$) =
  (1.000,2.76$\times10^{-3}$ eV$^2$/c$^4$),
  (0.937,2.50$\times10^{-3}$ eV$^2$/c$^4$) and
  (1.000,2.16$\times10^{-3}$ eV$^2$/c$^4$).

\end{itemize}





